% CV / Résumé — clean, ATS-friendly, single page, placeholder-driven.
% Tectonic is XeTeX: sources are always UTF-8, so no inputenc package.
\documentclass[11pt,a4paper]{article}

\usepackage[{{language}}]{babel}
\usepackage[T1]{fontenc}
\usepackage{textcomp}
\usepackage{lmodern}
\usepackage[a4paper,ignoreheadfoot,nomarginpar,top=15mm,bottom=15mm,left=15mm,right=15mm,headheight=12pt,headsep=4mm,footskip=8mm]{geometry}
\usepackage{array}
\usepackage{booktabs}
\usepackage{tabularx}
\usepackage{enumitem}
\usepackage{xcolor}
\usepackage{colortbl}
\usepackage{titlesec}
\usepackage{fancyhdr}
\usepackage{hyperref}

% Black links; PDF metadata comes from the placeholders.
\hypersetup{
  colorlinks=true,
  linkcolor=black,
  urlcolor=black,
  citecolor=black,
  filecolor=black,
  pdftitle={{{name}}},
  pdfauthor={{{name}}},
  pdfkeywords={{{headline}}}
}

% Greys from the reference CV: black!20 rules, grey secondary text.
\definecolor{cvgray}{gray}{0.55}
\definecolor{cvrule}{gray}{0.8}

% Empty header and footer: a one-page CV needs neither.
\pagestyle{fancy}
\fancyhf{}
\renewcommand{\headrulewidth}{0pt}

\setlength{\parindent}{0pt}
\setlength{\parskip}{2pt}

% Section titles: bold, a touch larger than the body, thin grey rule hugging
% the title (negative vspace pulls the rule up right under the text).
\titleformat{\section}{\normalfont\large\bfseries}{}{0pt}{}[\vspace{-8pt}\textcolor{cvrule}{\rule{\textwidth}{0.4pt}}]
\titlespacing*{\section}{0pt}{6pt}{3pt}

% Tight bullets, used under every entry.
\setlist[itemize]{noitemsep,topsep=1pt,leftmargin=1.2em,label=\textcolor{cvgray}{\textbullet}}

% Reusable entry: \cventry{role}{organization}{dates}{location}
% Role and organization share the left cell, dates and location the right one,
% so the extracted text reads role, organization, dates, location.
\newcommand{\cventry}[4]{%
  \begin{tabularx}{\textwidth}{@{}>{\raggedright\arraybackslash}X>{\raggedleft\arraybackslash}p{4.2cm}@{}}
    \textbf{#1}\newline #2 & \textit{#3}\newline #4 \\
  \end{tabularx}\vspace{-2pt}%
}

\begin{document}

  % ----------------------------------------------------------------------
  % Header — replace every {{...}} placeholder.
  % ----------------------------------------------------------------------
  \begin{center}
    {\huge\bfseries {{name}}}\\[4pt]
    {\large {{headline}}}\\[2pt]
    {\small
      \href{mailto:{{email}}}{{{email}}} \textbar\
      {{phone}} \textbar\
      {{location}} \textbar\
      \url{{{website}}}}
  \end{center}
  \vspace{-2pt}

  % ----------------------------------------------------------------------
  % Profile — one short paragraph, three to four lines.
  % ----------------------------------------------------------------------
  \section*{Profile}
  Generic software engineer with several years building reliable web
  applications and automation tools. Comfortable from data models and APIs to
  release pipelines, with a focus on clarity, testing and measurable results.

  % ----------------------------------------------------------------------
  % Experience — repeat the block for each role, newest first. \cventry takes
  % {role}{organization}{dates}{location}; add two to four bullets after it.
  % ----------------------------------------------------------------------
  \section*{Experience}

  \cventry{Senior Software Engineer}{Example Corp}{2022 -- Present}{City, Country}
  \begin{itemize}
    \item Led delivery of an internal service used by several teams, cutting a
    recurring manual task from hours to minutes.
    \item Built and documented the continuous integration pipeline for the main
    product.
    \item Mentored new engineers and reviewed design proposals.
  \end{itemize}

  \cventry{Software Engineer}{Sample Industries}{2020 -- 2022}{City, Country}
  \begin{itemize}
    \item Implemented and maintained customer-facing APIs and their database
    layer.
    \item Automated release checks, reducing failed deployments.
  \end{itemize}

  \cventry{Junior Developer}{Another Company}{2018 -- 2020}{City, Country}
  \begin{itemize}
    \item Built internal tools and dashboards that replaced manual spreadsheets.
  \end{itemize}

  % ----------------------------------------------------------------------
  % Education — two entries, same \cventry{degree}{institution}{dates}{city}.
  % ----------------------------------------------------------------------
  \section*{Education}

  \cventry{M.Sc. in Software Engineering}{Sample University}{2020 -- 2022}{City, Country}
  \cventry{B.Sc. in Computer Science}{Sample University}{2016 -- 2020}{City, Country}

  % ----------------------------------------------------------------------
  % Skills — two columns: category and a comma-separated list. Add rows as
  % needed; keep each list on one line if you can.
  % ----------------------------------------------------------------------
  \section*{Skills}
  {\renewcommand{\arraystretch}{1.2}%
    \arrayrulecolor{cvrule}%
    % No \toprule: the grey rule under the section title already caps the
    % table, and a second line right below it would read as a double rule.
    \begin{tabularx}{\textwidth}{@{}>{\bfseries}p{3.2cm}>{\raggedright\arraybackslash}X@{}}
      Languages & Python, JavaScript, SQL, Bash \\
      \midrule
      Frameworks & FastAPI, Node.js, Spring \\
      \midrule
      Tools & Git, Docker, Linux, CI/CD \\
      \bottomrule
    \end{tabularx}}

  % ----------------------------------------------------------------------
  % Projects — two entries, each with a one-line description and a link.
  % ----------------------------------------------------------------------
  \section*{Projects}

  \textbf{Project One} --- a small tool that automates a common task and is
  documented for reuse.\\
  \url{https://example.org/project-one}

  \vspace{2pt}
  \textbf{Project Two} --- a library used by other projects, with tests and
  examples.\\
  \url{https://example.org/project-two}

  % ----------------------------------------------------------------------
  % Languages — replace with your own levels.
  % ----------------------------------------------------------------------
  \section*{Languages}
  English (native) \textbar\ Spanish (professional)

\end{document}
